\documentclass[10pt]{article}

\usepackage[utf8]{inputenc}

\usepackage[T1]{fontenc}

\usepackage{amsmath}

\usepackage{amsfonts}

\usepackage{amssymb}

\usepackage[version=4]{mhchem}

\usepackage{stmaryrd}

\usepackage{graphicx}

\usepackage[export]{adjustbox}

\graphicspath{ {./images/} }



\begin{document}

при стремлении $z \rightarrow \zeta$ по любому некасательному пути; определение точки Плеснера $\zeta \in \Gamma \cdot$ нужно изменить так, чтобы рассматривались углы $\Delta$ с веріпиной $\zeta$ и сторонами, образующими углы, меньшие $\pi / 2$, с нормалью к $\Gamma$ в точке $\zeta$ (cм. [2]).



Аналогом П. т. в терминах категории множеств является Мейера теорема.



Лит.: [1] П л е с н е p А. И., “Успехи матем. наук", 1967, ства аналитических фунњций, 2 изд., М. П. Ј., 1950; [3] Л ов а т е р Д ж., в кн.: Итоги науки и техники. Математический



\includegraphics[max width=\textwidth, center]{2023_07_08_b3a169001b5c75bfc81eg-1}

КАЯ КРИВАЯ - множество точек $L$ действительной аффинной плоскости, координаты к-рых удовлетворяют уравнению



\[

f(x, y)=0,

\]



где $f(x, y)$ - многочлен степени $n$ от координат $x, y$; число $n$ наз. п о р я д к о м кривой $L$. Если многочлен $f$ приводим, т. е. разлагается на множители $f_{1}, \ldots, f_{k}$, то определяемая уравнением (1) кривая $L$ наз. п р ив о д и о й и я вляется объединением кривых $L_{1}$, ..., $L_{k}$ (к о м п о н е н т $L$ ), задаваемых соответственно уравнениями



\[

f_{1}=0, \ldots, f_{n}=0 .

\]



Если же многочлен $f$ неприводим, то кривая $L$ наз. в е п р и в о д и м о й. Две неприводимые П. д. а. к., одна из к-рых имеет порядок $n$, а другая - порядок $m$, пересекаются не более чем в $m n$ точках (т е о р е ма Б $\boldsymbol{\text { в }} \mathbf{\text { y). }}$



Одна и та же П. д. а. к. $L$ может определяться различными уравнениями. Пусть $I_{L}$ - совокупность многочленов, обращающихся в нуль во всех тонках кривой $L$. Если $L$ неприводима, то из того, что $f g=0$ на $L$, следует равенство нулю $f$ или $g$; в этом случае факторкольцо $K(L)=K / I_{L}$ (здесь $K$ - кольцо всех многючленов) не. имеет делителей нуля и наз. к о л ь ц о м м н о го ч л е н о в на $L$.



С неприводимой П. д. а. к. $L$ ассоциируется также нек-рое поле $K(L)$, наз. п о л е м р а ц и о н а лі ь ны ф у н к ц и й на $L$. Оно состоит из рациональных фуункций $\frac{p(x, y)}{q(x, y)}$, причем $q$ не делится на $f$, pacсматриваемых с точностью до равенства на $L\left(\frac{p}{q}\right.$ и $\frac{\tilde{p}}{\tilde{q}}$ наз. равными на кривой $L$, заданной уравнением (1), если многочлен $p \tilde{q}-\tilde{p} q$ делится на $f)$. Поле $K(L)$ является полем частных кольца $K(L)$.



Отображение $F:(x, y) \rightarrow \varphi(x, y), \psi(x, y)$ плоскости в себя наз. р е г у л я р ны м на П. д. а. к. $L$, если $\varphi, \psi \in K(L)$. Кривые $L$ и $M$ наз, и з о м о р ф н ы м и, если существуют регулярные (соответственно на $L$ и на $M$ ) отображения $F: L \rightarrow M$ и $G: M \rightarrow L$, обратные друг другу; при этом кольца $K(L)$ и $K(M)$ ияоморфны. В частности, изоморфны аффинно эквивалентные кривые.



Более общим является р а ц и о в а ль ь о е о т об р а ж е и е кривой $L$ в кривую $M$, осуществляемое рациональными функциями. Оно устанавливает соответствие между всеми точками кривых, кроме конечного их числа, и определяется так. Пусть $f=0$ и $g=0-$ уравнения кривых $L$ и $M$ соответственно, тогда рациональное отображение $F$ задается такой парой рациональных функций $\varphi$ и $\psi$, определенных на $L$, что $g(\varphi$, $\psi)=0$ на $M$. Кривые $L$ и $M$ наз. 6 и р а ц и о на л ьн о э к в и в а л е н т ным и, если существуют рацұональные отображения $L$ в $M$ и $M$ в $L$, обратные друг другу; при этом поля $k(L)$ и $k(M)$ оказываются изоморфными. Такие рациональные отображения наз. бирац и о нальны ми, или кремоновы ми, преобразованиями. Все кремоновы преобразования на плоскости могут быть осуществлены последовательным проведением преобразований вида $x \rightarrow \frac{1}{x}, y \rightarrow \frac{1}{y}$. Бирациональная әквивалентность - более грубое отношение, чем изоморфизм, однако классификация П. д. а. к. с этой точкй зрения оказывается проще и обозримее.



Простейшим примером рационального отображения является проективное преобразование. Важную роль играет д у а л ь н о е о т о б р а же н и е неприводимой кривой $L$, отличной от прямай, в кривую $L^{*}$, д уа л ь н у ю $L$, определяемое формулами:



\[

u=\frac{\frac{\partial f}{\partial x}}{f-x \frac{\partial f}{\partial x}-y \frac{\partial f}{\partial y}}, \quad v=\frac{\frac{\partial f}{\partial y}}{f-x \frac{\partial f}{\partial x}-y \frac{\partial f}{\partial y}},

\]



где $f$ - многочлен, определяющий $L$. Уравнение



\[

g(u, v)=0

\]



определяющее кривую $L *$, получается исключением $x, y$ из уравнений (1), (2). Вследствие связи дуального отображения с касательным преобразованием саму кривую $L^{*}$ в ряде случаев можно предетавлять как огибающую семейства прямых, касательных к $L$.



Порядок $L^{*}$ наз. к л а с с о м $n^{*}$ кривой $L$. Отношение дуальности взаимно, т. е. $L^{* *}=L$, и является отражением двойственности принципа проективњой геометрии.



Точка $x$ П. д. а. к. $L$, заданной уравнением (1), наз. о с о б о й т оч к о й, если в ней $\operatorname{grad} f=0$. Анализ особенностей является необходимым элементом иеследования кривой $L$, однако полная классификация особенностей өще далека от завершения (1983).



Если все производные многочлена $f$ до $(r-1)$-го порядка включительно в точке $x$ равны нулю, а хотя бы одна производная $r$-го порядка отлична от нуля в $x$, то $x$ наз. т оч к о й к р а т н о с т и $r$ и притом обыкновенной, если в ней существуют $r$ различных касательных. Примеры особых точек:



\begin{enumerate}

  \item $x^{3}-x^{2}+y^{2}=0 ;(0,0)$ - обыкновенная двойная точка, точка самопересечения;



  \item $x^{2}+x^{3}+y^{2}=0 ;(0,0)$ - изолированная точка;



  \item $x^{3}+y^{2}=0 ;(0,0)-$ точка заострения, возврата;



  \item $2 x^{4}-3 x^{2} y+y^{2}-2 y^{3}+y^{4}=0 ;(0,0)$ - точка самоприкосновения.



\end{enumerate}



Неособая точка $x$ П. д. а. к. $L$, заданной уравнением (1), наз. т о ч к о й п е р е г и б а, если в ней



\[

H(x, y)=\left|\begin{array}{lll}

\frac{\partial^{2} f}{\partial x^{2}} & \frac{\partial^{2} f}{\partial x \partial y} & \frac{\partial f}{\partial x} \\

\frac{\partial^{2} f}{\partial x \partial y} & \frac{\partial^{2} f}{\partial y^{2}} & \frac{\partial f}{\partial y} \\

\frac{\partial f}{\partial x} & \frac{\partial f}{\partial y} & 0

\end{array}\right|=0 .

\]



Другими словами, точки перегиба - это точки пересечения кривой $L$ и кривой $H$, задаваемой уравнением (3) и называемой г е c c и а в о й кривой $L$. Точкам перегиба кривой $L$ соответствуют на дуальной кривой $L^{*}$ точки возврата.



Для всякой П. Д. а. к. имеет место соотношение (Ф. Клейн, F. Klein, 1876):



\[

n+2 d+r=n^{*}+2 d^{*}+r^{*} \text {, }

\]



где $n$ - порядок $L, n^{*}-$ класс $L, r^{*}$ - qисло точек перегиба $L, d^{*}$ - число изолированных двойных касательных к $L$ (двойных точек $L^{*}$ ), $r$ - число точек возврата $L$ (точек перегиба $\left.L^{*}\right), d$ - число двойных точек $L$. См. также Плюккера боржулы.



Любая неприводимая плоская кривая $L$ бирационально эквивалентна неприводимой кривой $L_{0}$, имеющей лишь обыкновенные особенности.





\end{document}